% ============================================================
% 附录 C：概念批注 —— 与 nLab 的对照
% 本文件由 Lead 撰写；nLab 条目名均已实测（2026 年 11 月检索）。
% 收录标准：原文中反复出现的概念，且 nLab 确有独立条目或显著段落。
% 来源分次：nLab 为 B 级，一律与 A 级文献（原文、Bondal–Kapranov、
% Bondal–Orlov、Huybrechts 等）交叉；未查到独立条目者如实注明。
% ============================================================
\chapter{概念批注：与 nLab 的对照}
\label{app:nlab}

本书正文是 2003 年的文字，术语与论证风格属于 Gabriel--Zisman 分式演算与
Verdier 三角范畴那一代。二十年来，同一批对象在 nLab 上被重新组织过一遍：
表述换成了 $(\infty,1)$-范畴与稳定同伦论的语言，重点也从\emph{公理怎么列}
移到\emph{这些公理记的是哪个更高结构的影子}。本附录挑出原文中反复出现的
概念，逐条指出 nLab 对应条目在讲什么、与 Orlov 的用法差在哪里、以及读到
正文哪一处时值得回头对照。写成这种形式有实际用处：正文里那些看起来只是
技术性铺垫的引理（例如 Serre 函子的唯一性与正合性），在 nLab 的语言里
往往正是被单独立为命题的那几条，读者由此可以判断哪些地方可以快读、
哪些地方必须逐字核对。

下面每条的“出处”一栏分两级标注：A 级为原始论文或经典教材，B 级为 nLab
条目（含条目名，便于检索）。凡 nLab 的说法与 A 级文献不完全一致处，均已注明；
凡 nLab 没有独立条目、只能从相关条目里间接读到的，也如实写明。

\section{用法说明}

本附录不替代正文，也不替代 nLab：它只回答“这个概念在另一套语言里叫什么、
在哪里找、彼此差多少”三个问题。读法建议如下。

\begin{enumerate}
  \item 先读正文，遇到卡住的概念再查本附录，不要反过来用本附录去预习正文 ——
        nLab 的组织方式（以 $(\infty,1)$-范畴为骨架）与本文的次序并不吻合，
        顺着它读会得到一张漂亮但不合用的地图。
  \item 凡本附录说“差异在于……”的地方，差异都是\emph{表述层面}的：
        两边定义的数学内容相同，只是要求的附加条件有多有少。这类差异在
        引用时会造成实质区别（例如半正交分解是否要求成员可容许），故不可略过。
  \item 本附录的每条批注都可以独立阅读，不依赖前后条。
\end{enumerate}

\section{逐条对照}

\begin{conceptlist}

\concept{三角范畴}{triangulated category}%
{原文第 1.1 节把三角范畴作为\emph{原始对象}给出：一个带平移函子的加法范畴，
配一族特异三角，满足公理 T1--T4（其中 T4 即八面体公理）。nLab 的同名条目
采取另一种次序：先有稳定 $(\infty,1)$-范畴，三角范畴是其同伦范畴，
三角结构本身则是“不可函子化的影子”—— 映射锥依赖选择，故三角范畴之间的
函子不自动保锥。条目另设“enhanced triangulated category”专讲这件事。
差异的实用后果是：本文在 3.1 节用 Postnikov 系统逐层粘合映射锥时，
粘合的结果在同构意义下才唯一；nLab 的框架里这一步是取 $(\infty,1)$-极限，
不存在选择问题。两种写法给出同一个对象，但前者需要读者自己验证无关性。}%
{nLab “triangulated category”、“enhanced triangulated category”（B 级）；
Verdier 1996、Neeman 2001（A 级）。}

\concept{特异三角}{distinguished triangle}%
{原文的记法是 $X\to Y\to Z\to X[1]$，并称之为特异三角；nLab 条目名为
“distinguished triangles”，内容与原文一致，只是把它放在稳定 $(\infty,1)$-范畴
的纤维序列（fiber sequence）这一更大类里讲。读 1.1 节的公理 T2、T3 时，
可以拿 nLab 里“旋转”与“同构封闭”的表述互校，两边逐字相同。}%
{nLab “distinguished triangles”（B 级）；原文公理 T1--T4（A 级）。}

\concept{正合函子}{exact functor}%
{原文定义 1.1.1 要求函子把特异三角送到特异三角；nLab 的 “exact functor”
条目在三角范畴的语境下给出同样的定义，并额外说明：在稳定 $(\infty,1)$-范畴
层面，“正合”等价于“保持有限极限”，从而三角范畴层面的正合性只是这一条的
影子。本文第 2 章用正合性做的所有论证（复合、伴随、张量积）都在三角范畴
层面完成，不需要动到 $(\infty,1)$ 层。}%
{nLab “exact functor”（B 级）；原文定义 1.1.1（A 级）。}

\concept{导出范畴}{derived category}%
{原文定义 1.2.3 走的是局部化路线：取同伦范畴 $H(\cA)$，再关于零同伦复形
构成的三角子范畴作局部化，得到 $\Du(\cA)$。nLab 的同名条目则从链复形范畴
出发，在拟同构处作 $(\infty,1)$-局部化，再取同伦范畴；两条路线给出同一个
三角范畴（这一点是经典结果）。差异在于谁承担“高阶信息”：原文把它交给
分式演算的 Ore 条件，nLab 把它交给局部化的更高结构。本文几乎不显式使用
高阶信息，故两套语言在此处的差别对阅读影响最小。}%
{nLab “derived category”、“stable (∞,1)-category”（B 级）；
Gabriel--Zisman 1967、Verdier 1996（A 级）。}

\concept{Serre 函子}{Serre functor}%
{这是全篇最值得对照的一条。原文定义 1.1.3 给出 Serre 函子，引理 1.1.4 证明
它在分次自然同构意义下唯一，引理 1.1.5 证明它在三角范畴中必正合。nLab 的
“Serre functor” 条目做了三件正好对应的事：其一，同时列出 Bondal--Kapranov
原始定义与后世较弱的定义，并指出原定义中“与 $S$ 诱导的同构交换”这一条
其实可由 $\phi$ 的自然性推出，故可删；其二，把“任何自等价与 Serre 函子
交换”列为命题；其三，把“三角范畴中的 Serre 函子必正合”与“任意两个
Serre 函子由分次自然同构相连”分别列为命题 —— 后两条恰是引理 1.1.4 与
引理 1.1.5。换言之，本文当作预备知识顺手引用的两条引理，在 nLab 那里是
条目的主结论。条目还给出光滑射影簇上 $(\cdot\otimes\omega_X)[n]$ 这个例子，
与本文 1.3 节的用法完全一致。}%
{nLab “Serre functor”（B 级）；Bondal--Kapranov 1989（A 级）；
原文定义 1.1.3、引理 1.1.4、引理 1.1.5（A 级）。}

\concept{完美复形}{perfect complex}%
{原文 1.3 节用完美复形来陈述导出范畴里的有限性条件。nLab 的条目名是
“perfect chain complex”（以 “perfect complex” 访问会被重定向到该页），
定义与原文相同：在环化空间上，局部拟同构于一个有界、每项为有限秩自由
$\cO_X$-模的复形。值得记下的是该条目强调的两点：其一，在 $(\infty,1)$ 语言里
完美复形就是可对偶对象（dualizable object）；其二，“完美”与“紧”并不
总是同一回事，只有在拟紧拟分离概形上才由 Thomason--Trobaugh、Neeman、
Bondal--van den Bergh 的工作保证一致。本文始终在光滑射影簇上工作，
两个概念重合，故原文可以不加区分地使用。}%
{nLab “perfect chain complex”（B 级）；Thomason--Trobaugh 1990、
Bondal--van den Bergh 2003（A 级）。}

\concept{可容许子范畴}{admissible subcategory}%
{原文 2.2 节开头定义：三角范畴 $\cD$ 的全三角子范畴 $\cN$ 若含入函子有右伴随
（左伴随），则称右（左）可容许；两者兼备才称可容许。nLab 的
“admissible subcategory” 条目逐字给出同样的三层区分，并补一句本文未写、
但恰好解释了本文为何用得如此顺手的话：若 $\cD$ 有 Serre 函子，则右可容许
与左可容许等价。原文在光滑射影簇的导出范畴上工作，Serre 函子总存在
（命题 1.3.1），故原文只须验证一侧。}%
{nLab “admissible subcategory”（B 级，另见 Stacks Project 标签 13.40）；
原文 2.2 节定义（A 级）。}

\concept{半正交分解}{semiorthogonal decomposition}%
{这里有一处必须留意的实质差异。nLab 的 “semiorthogonal decomposition” 条目
只列两条公理：$i<j$ 时 $\Hom(\cA_j,\cA_i)=0$；各 $\cA_i$ 共同生成 $\cD$。
原文定义 2.2.3 则明确要求序列 $(\cN_0,\dots,\cN_n)$ 的每一项都是
\emph{可容许}子范畴。在带 Serre 函子的三角范畴里两种写法等价（由前一条批注），
故对本文的定理没有影响；但若把原文的结论搬到一般三角范畴上引用，
“可容许”这一条不能省，否则 $\langle\cA_1,\dots,\cA_n\rangle$ 生成出来的分解未必能
拼回原范畴。阅读 2.2 节时，建议先按原文的定义记住“可容许”这个条件，
再回看 nLab 只列两条公理的版本，两者的落差正好提示了 Serre 函子在这里的作用。}%
{nLab “semiorthogonal decomposition”（B 级）；原文定义 2.2.2、2.2.3（A 级）。}

\concept{Fourier--Mukai 变换}{Fourier--Mukai transform}%
{原文式 (1) 定义 $\Phi_{\cE}(\cdot):=\Rb q_*(\cE\otL p^*(\cdot))$，全文围绕
“哪些正合函子可写成这种形状”展开。nLab 的 “Fourier--Mukai transform” 条目
与本文关系最直接：它把积分核的定义写成
$F\mapsto \Rf q_*(\Lf p^*(F)\otL \cE)$，与式 (1) 逐字一致；随后以
“定理（Orlov）”的形式陈述本文定理 3.2.1，并明确注明证明见
“Orlov 2003, 3.2.1” —— 也就是说，本译本第三章的主定理正是 nLab 该条目
所引的那一条。条目另有三点本文未及展开、但读者值得知道的补充：
（一）原文定理 3.2.1 假设函子有右伴随，而 Bondal--van den Bergh 2002 证明
三角全忠实函子自动有右伴随，故该假设可去；
（二）曾有人猜测所有三角函子都是 Fourier--Mukai 型，Rizzardo--Van den Bergh
2015 给出反例，此路不通；
（三）在 $(\infty,1)$ 增强层面，$\mathrm{Func}^L(\Du(X),\Du(Y))\simeq \Du(X\times Y)$，
这是定理 3.2.1 的精确化版本。}%
{nLab “Fourier-Mukai transform”（B 级，条目内引 Orlov 2003 §3.2.1、
Bondal--van den Bergh 2002、Rizzardo--Van den Bergh 2015）；
Mukai 1981、原文定理 3.2.1（A 级）。}

\concept{Postnikov 系统}{Postnikov system}%
{须提醒同名不同用。原文 3.1 节的 Postnikov 系统指的是三角范畴中的一列对象
$X_0\to X_1\to\cdots$，相邻两级由特异三角相连，再用卷积 $\Tot$ 把这些层
拼成一个对象；它借的是拓扑里 Postnikov 塔的\emph{意象}。nLab 的
“Postnikov system” 条目讲的则是链复形的 $n$-像分解（把链映射分解为
$\im_{n+1}(f)_\bullet$ 的过滤）以及单纯预层的 Postnikov 塔，是拓扑那一侧的
严格版本。两者不是同一个定义，只有“逐层截断再拼回”这一条精神相通。
读 3.1 节时不要拿 nLab 该条目替代原文定义。}%
{nLab “Postnikov system”（B 级，内容为链复形 $n$-像分解与单纯预层）；
原文 3.1 节（A 级）。}

\concept{Bondal--Orlov 重构定理}{Bondal-Orlov reconstruction theorem}%
{原文定理 2.1.3 是它的一个特例：由 $\Db(X)$ 作为三角范畴的等价类还原出
光滑射影簇 $X$ 本身（要求典范层或反典范层丰富）。nLab 设有同名条目，
把它作为独立定理陈述，并给出 Serre 函子与典范层如何被范畴结构读出的证明
思路 —— 这条思路比原文的证法更贴近现代叙述，可作为 2.1 节的补充读物。}%
{nLab “Bondal-Orlov reconstruction theorem”（B 级）；
Bondal--Orlov 2001、原文定理 2.1.3（A 级）。}

\concept{例外族}{exceptional collection}%
{原文定义 2.2.4 给出例外对象与完备例外族，例 2.2.5 给出 $\PP^N$ 上的
$(\cO(i),\cO(i+1),\dots,\cO(i+N))$。就本次检索所见，nLab 没有独立的
“exceptional collection” 条目：这一概念只在 “semiorthogonal decomposition”
条目里以“例子的主要来源”的形式被提到一句。因此本条的出处只用 A 级文献，
不引 nLab，以免读者按条目名找不到页面。}%
{Bondal 1989、Rudakov 1990、原文定义 2.2.4 与例 2.2.5（A 级）；
nLab 侧仅 “semiorthogonal decomposition” 条目提及，无独立条目，故不引。}

\end{conceptlist}

\clearpage
\section{一条总的读法}

把上面十二条放在一起看，可以看出两种语言的分工：原文把每个概念当作
\emph{可用的工具}，定义了就用，用完就走；nLab 则习惯先把工具放回它所属的
更高结构里，说明它是那个结构投下的影子，再回到具体计算。前者适合把一篇
90 页的综述从头读到尾，后者适合在某一处卡住时追问“这个条件为什么可以省”。
两者不冲突，也不互相替代。

对本书读者，一个实用的建议是：第一次读时完全按原文的次序走，把
$(\infty,1)$ 语言暂时搁置；第二次读时，把本附录当作注释反过来查 ——
凡正文中“显然”“不难验证”之处，多半正是 nLab 条目里被单独立为命题的地方；
凡正文中条件写得比预期更严之处（如半正交分解要求可容许），多半正是
两套语言在这一点上分道的地方。
